\documentclass[11pt]{article}
\usepackage[margin=1in]{geometry}
\usepackage{amsmath,amssymb,amsthm}
\usepackage[T1]{fontenc}
\usepackage{lmodern}
\usepackage{graphicx}
\usepackage{hyperref}
\usepackage{enumitem}

\newtheorem{theorem}{Theorem}
\newtheorem{lemma}[theorem]{Lemma}
\newtheorem{corollary}[theorem]{Corollary}
\newtheorem{remark}[theorem]{Remark}

\newcommand{\ellT}{\ell_2}
\newcommand{\cinf}{\ell_\infty}
\newcommand{\czz}{c_{00}}
\newcommand{\TT}{\mathbb T}
\newcommand{\supp}{\operatorname{supp}}
\newcommand{\Span}{\operatorname{span}}
\newcommand{\Dop}[1]{D_{#1}}

\title{Self-similar twisted Hilbert extensions are uniformly block-natural centralizer extensions}
\author{Codex agent lane 9}
\date{July 22, 2026}

\begin{document}
\maketitle

\section*{Status and source conjecture}

This is a candidate partial result, subject to human review.  Corr\^ea,
Dantas, and Rodr\'iguez-Vidanes prove that the Kalton--Peck space is the
unique nontrivial self-similar twisted Hilbert space among the spaces
$Z(\varphi)$ generated by a scalar Lipschitz function, and then conjecture
that the same is true for every twisted Hilbert space
\cite[Theorem 7.4 and p.~19]{CDR}:

\begin{center}
  \includegraphics[width=0.96\linewidth]{figures/open_problem_crop.png}
\end{center}

The exact conjecture is:
\begin{quote}
The Kalton--Peck space $Z_2$ is, up to isomorphism, the only nontrivial
self-similar twisted Hilbert space.
\end{quote}

The theorem below does not finish the classification.  It reduces the full
problem from arbitrary quasilinear extensions of $\ell_2$ by $\ell_2$ to
extensions generated by exact $\ell_\infty$-centralizers.  Moreover, the
canonical block defects of the resulting centralizer have one bound valid for
\emph{all} normalized block bases.  The associated scalar indicator therefore
satisfies a dimension-free approximate Shannon chain rule for every ordered
refinement of a finite probability vector.  In particular, a counterexample
cannot arise solely from a genuinely non-centralizer quasilinear map, and the
remaining classification has a precise scalar cocycle obstruction.

\section*{Definitions}

We work over $\mathbb R$ or $\mathbb C$.  Let
$G=\{-1,1\}^{\mathbb N}$ be the Cantor sign group in either scalar field.
For $u\in G$, $\Dop u$ denotes the diagonal isometry
$(x_n)\mapsto(u_nx_n)$ of $\ell_2$.

A homogeneous map $\Omega:\czz\to\ell_2$ is \emph{quasilinear} if there is
$Q<\infty$ such that
\[
 \|\Omega(x+y)-\Omega(x)-\Omega(y)\|_2
 \le Q(\|x\|_2+\|y\|_2) \qquad (x,y\in\czz).
\]
It represents the exact sequence whose middle space is the completion of
$\ell_2\times\czz$ for the quasinorm
\[
 \|(y,x)\|_\Omega=\|y-\Omega(x)\|_2+\|x\|_2.
\]
Replacing $\Omega$ by $\Omega+L+B$, with $L:\czz\to\ell_2$ linear and
$B:\czz\to\ell_2$ bounded homogeneous, gives an equivalent exact sequence.

The map $\Omega$ is \emph{sign-exact} if
\[
 \Omega(\Dop u x)=\Dop u\Omega(x) \qquad (u\in G,x\in\czz).
\]
In the complex case, \emph{exact} means the same identity for every
$u\in\TT^{\mathbb N}$; over the reals, sign-exact and exact coincide.
It is an $\ell_\infty$-\emph{centralizer} if there is $C<\infty$ such that
\[
 \|\Omega(ax)-a\Omega(x)\|_2
 \le C\|a\|_\infty\|x\|_2
 \qquad(a\in\ell_\infty,x\in\czz).
\]
This is the standard sequence-space centralizer estimate; see, for example,
\cite[Sec.~3]{CG}.

For a normalized block basis $(u_n)$ of the canonical basis of $\ell_2$,
write $U=[u_n]$ and let $T_U:\ell_2\to U$ be the isometry
$T_Ue_n=u_n$.  The extension $X=X_\Omega$ is self-similar with respect to
$(u_n)$ if there is an isomorphism from the original exact sequence to its
block restriction inducing $T_U$ on both the kernel and quotient.  It is
self-similar if this holds for every normalized block basis, as in
\cite[Sec.~7]{CDR}.

\section*{Idea of the proof}

Singleton block bases with arbitrary coordinate signs show that the
extension class of $\Omega$ is fixed by every sign-diagonal conjugacy.  This is
initially only an algebraic assertion: for every sign sequence there is a
possibly different, possibly unbounded linear correction and a bounded error
whose bound may depend on the sign sequence.  The first Baire argument makes
the bounded errors uniform.  An invariant mean on the compact sign group
can then be applied safely on each finite-support vector, producing an
equivalent sign-exact representative $\omega$.

Sign-exactness makes $\omega$ support-preserving.  To test multiplication by an
arbitrary $a\in B_{\ell_\infty}$, encode $a_n$ as the first coordinate of a
unit vector in a disjoint two-dimensional block.  Self-similarity for this
block basis turns the multiplier commutator $a\omega-\omega a$ into a bounded
map, up to a bounded diagonal correction.  Thus each individual multiplier
has a finite commutator bound.  The second Baire argument promotes these
pointwise finite bounds to one uniform bound over the unit ball of
$\ell_\infty$, which is exactly the centralizer condition.  In the complex
case, the centralizer bound then permits a final bounded average over
$\TT^{\mathbb N}$, giving a fully exact representative.

There is a further uniformity mechanism.  If the commutator constants for
increasing coordinate injections were unbounded, choose successively worse
finite witnesses, move their domain and range coordinates into disjoint
windows, and paste all of them into one increasing injection.  The resulting
single commutator would be unbounded, contrary to self-similarity.  The same
gliding-hump paste, now using the uniform injection bound on both the domain
and range, gives one constant for the canonical defects of every normalized
block basis.  Pairing that vector-valued defect with the corresponding block
sum turns it into the approximate chain rule stated below.

\section*{The reduction theorem}

\begin{theorem}\label{thm:main}
Let
\[
 0\longrightarrow\ell_2\longrightarrow X\longrightarrow\ell_2
 \longrightarrow0
\]
be a self-similar twisted Hilbert extension.  Then this exact sequence is
equivalent to one generated by a quasilinear map
$\omega:\czz\to\ell_2$ satisfying
\begin{enumerate}[label=\textup{(\roman*)}]
\item $\omega(e_n)=0$ for every $n$;
\item $\omega(\Dop u x)=\Dop u\omega(x)$ for every unimodular diagonal
      $u$ (all $u\in\TT^{\mathbb N}$ over $\mathbb C$, and all signs over
      $\mathbb R$) and $x\in\czz$;
\item there is $C<\infty$ such that
\[
 \|\omega(ax)-a\omega(x)\|_2
 \le C\|a\|_\infty\|x\|_2
 \qquad(a\in\ell_\infty,x\in\czz).
\]
\end{enumerate}
Thus every self-similar twisted Hilbert extension is centralizer-generated.
\end{theorem}

The proof also gives the following strengthening.  It is stated separately so
that the scalar obstruction left by the source conjecture is explicit.

\begin{theorem}[Uniform block naturality and indicator chain rule]
\label{thm:uniform-block}
The representative $\omega$ in Theorem~\ref{thm:main} can be chosen so that
there is a constant $K<\infty$, independent of the normalized block basis
$U=(u_n)$, satisfying
\begin{equation}\label{eq:uniform-canonical}
 \left\|T_U\omega(x)-\omega(T_Ux)+
       \sum_nx_n\omega(u_n)\right\|_2
 \le K\|x\|_2 \qquad(x\in\czz).
\end{equation}

For every finite positive vector $p=(p_n)$ put
\[
 \mathcal H_\omega(p)=
 \left\langle\omega\big((\sqrt {p_n})_n\big),
                     (\sqrt {p_n})_n\right\rangle .
\]
If $q^{(n)}=(q_{n,k})_k$ are finite probability vectors and
$r=(p_nq_{n,k})_{n,k}$ is their ordered concatenated refinement, then
\begin{equation}\label{eq:entropy-chain}
 \left|\mathcal H_\omega(r)-\mathcal H_\omega(p)
       -\sum_np_n\mathcal H_\omega(q^{(n)})\right|
 \le K\sum_np_n .
\end{equation}
Here the finite vectors may be placed in arbitrary successive coordinate
blocks; changing their placement changes only the same uniform error class.
\end{theorem}

We divide the proof into lemmas.

\begin{lemma}[Block commutator criterion]\label{lem:blockcriterion}
Let $\Omega:\czz\to\ell_2$ represent the extension and let $(u_n)$ be a
normalized block basis.  If the extension is self-similar with respect to
$(u_n)$, then there are a linear map $L:\czz\to U$ and $C<\infty$ such that
\[
 \|T_U\Omega(x)-\Omega(T_Ux)+Lx\|_2\le C\|x\|_2
 \qquad(x\in\czz).
\]
\end{lemma}

\begin{proof}
Any operator making the defining exact-sequence diagram commute has, on the
dense algebraic model, the triangular form
\[
 T(y,x)=(T_Uy+Lx,T_Ux)
\]
for a linear $L:\czz\to U$.  Apply boundedness of $T$ to
$(\Omega(x),x)$, whose $\Omega$-quasinorm is $\|x\|_2$.  The first component
of the resulting target quasinorm is precisely the expression in the
statement.  This is the general part of the calculation used in
\cite[Lemma 7.1]{CDR}.
\end{proof}

First subtract the linear map determined by $e_n\mapsto\Omega(e_n)$.  We may
and do assume from now on that
\begin{equation}\label{eq:zero-basis}
 \Omega(e_n)=0 \qquad(n\in\mathbb N).
\end{equation}

\begin{lemma}[Uniform diagonal equivalence]\label{lem:uniform-diagonal}
There is $M<\infty$ such that, for every $u\in G$, one can find a linear
map $L_u:\czz\to\ell_2$ with
\begin{equation}\label{eq:uniform-diagonal}
 \|\Dop u\Omega(x)-\Omega(\Dop u x)+L_ux\|_2
 \le M\|x\|_2 \qquad(x\in\czz).
\end{equation}
\end{lemma}

\begin{proof}
For $m\in\mathbb N$, let $E_m$ be the set of $u\in G$ for which such an
$L_u$ exists with bound $m$.  The singleton block basis $(u_ne_n)$ and
Lemma~\ref{lem:blockcriterion} give $G=\bigcup_mE_m$.

We claim that $E_m$ is closed in the product topology of $G$.  Suppose
$u^{(j)}\to u$ and choose $L_j$ with error at most $m$.  By
\eqref{eq:zero-basis}, the diagonal commutator vanishes at every $e_n$;
hence $\|L_je_n\|_2\le m$.  A diagonal weak-subsequence argument in the
reflexive space $\ell_2$ gives vectors $v_n\in\ell_2$ such that
$L_je_n\to v_n$ weakly for every $n$.  Define the linear map $L$ on $\czz$
by $Le_n=v_n$.  For each fixed $x\in\czz$, product convergence in the sign
group implies $\Dop{u^{(j)}}x=\Dop ux$ eventually.  The remaining diagonal
action on the fixed vector $\Omega(\Dop ux)\in\ell_2$ converges strongly.
It follows that the diagonal commutator converges in norm.
Passing to the weak limit in the error inequality yields the same inequality
with $L$ and bound $m$.  Thus $u\in E_m$.

The sets satisfy
\[
 E_m^{-1}\subset E_m,
 \qquad E_mE_k\subset E_{m+k},
\]
because diagonal conjugacies are isometries and the corresponding linear
corrections may be conjugated and added.  Baire category gives an $m$ for
which $E_m$ has nonempty interior.  Therefore $E_m^{-1}E_m$ contains an
open neighborhood $W$ of the identity and $W\subset E_{2m}$.  The subgroup
generated by $W$ is open in the compact group $G$, hence has finite index;
compactness also gives an integer $r$ for which that subgroup is contained
in $W^r$.  Choose finitely many coset representatives and finite bounds for
them.  The product rule then gives one bound $M$ valid on all of $G$.
\end{proof}

\begin{lemma}[Sign exactification]\label{lem:exactification}
The map $\Omega$ is equivalent to a sign-exact quasilinear map
$\omega:\czz\to\ell_2$ satisfying $\omega(e_n)=0$.
\end{lemma}

\begin{proof}
Rewrite \eqref{eq:uniform-diagonal}, after multiplying by $\Dop u^{-1}$,
as
\[
 \Dop u^{-1}\Omega(\Dop u x)-\Omega(x)=A_u x+B_u(x),
\]
where $A_u$ is linear and $\|B_u(x)\|_2\le M\|x\|_2$.  The sign convention
is immaterial.  Since the left side vanishes at $e_n$, we may choose the
decomposition so that $\|A_ue_n\|_2\le M$.  Consequently, for every fixed
$x\in\czz$, both $u\mapsto A_ux$ and $u\mapsto B_u(x)$ are bounded
$\ell_2$-valued functions.

Let $\mathfrak m$ be an invariant mean on the compact abelian group $G$ and
define weakly in the reflexive space $\ell_2$
\[
 \omega(x)=\mathfrak m_u\big[\Dop u^{-1}\Omega(\Dop u x)\big].
\]
The uniform quasilinearity constant of the conjugates shows that $\omega$ is
quasilinear.  Invariance of the mean gives
\[
 \omega(\Dop v x)=\Dop v\omega(x) \qquad(v\in G).
\]
Moreover,
\[
 \omega-\Omega=\mathfrak m_u A_u+\mathfrak m_u B_u,
\]
the first summand being linear on $\czz$ and the second bounded by $M$.
Thus $\omega$ is equivalent to $\Omega$.  Finally $\omega(e_n)=0$ follows
from \eqref{eq:zero-basis}.
\end{proof}

Sign-exactness implies the support property
\begin{equation}\label{eq:support}
 \supp\omega(x)\subset\supp x.
\end{equation}
Indeed, choose a diagonal unit equal to $1$ on $\supp x$ and different from
$1$ off that support, and compare $\omega(ux)$ with $u\omega(x)$.

\begin{lemma}[Canonical block correction]\label{lem:canonical}
Let $(u_n)$ be a normalized block basis.  There is $C_U<\infty$ such that
\begin{equation}\label{eq:canonical}
 \left\|T_U\omega(x)-\omega(T_Ux)+
       \sum_nx_n\omega(u_n)\right\|_2
 \le C_U\|x\|_2
 \qquad(x\in\czz).
\end{equation}
\end{lemma}

\begin{proof}
By Lemma~\ref{lem:blockcriterion}, the first two terms become bounded after
adding a linear $L:\czz\to U$.  Average $L$ over the diagonal group on the
domain, acting on the range by multiplication of the whole $n$th block by
the same scalar.  Sign-exactness of $\omega$ shows that the averaged linear map
$\Lambda$ still gives a bounded correction and that $\Lambda e_n$ is
supported in the block supporting $u_n$.

At $e_n$, the uncorrected commutator equals $-\omega(u_n)$.  Hence
$\|\Lambda e_n-\omega(u_n)\|_2$ is bounded independently of $n$.  Because
these differences have disjoint supports, the linear map sending
$e_n$ to $\Lambda e_n-\omega(u_n)$ is bounded on $\ell_2$.  Replacing
$\Lambda e_n$ by $\omega(u_n)$ therefore changes only the final constant
and yields \eqref{eq:canonical}.
\end{proof}

\begin{lemma}[Uniform increasing-injection commutators]
\label{lem:uniform-injections}
For every increasing injection $\alpha:\mathbb N\to\mathbb N$, let
$J_\alpha e_n=e_{\alpha(n)}$.  There is a constant $D<\infty$, independent of
$\alpha$, such that
\begin{equation}\label{eq:uniform-injections}
 \|J_\alpha\omega(x)-\omega(J_\alpha x)\|_2
 \le D\|x\|_2 \qquad(x\in\czz).
\end{equation}
\end{lemma}

\begin{proof}
For a fixed $\alpha$, Lemma~\ref{lem:canonical} applied to the singleton
block basis $(e_{\alpha(n)})$ gives a finite bound because
$\omega(e_{\alpha(n)})=0$.  Denote the least such bound by $d(\alpha)$.

Suppose that $\sup_\alpha d(\alpha)=\infty$.  We construct one increasing
injection $\beta$ on successive finite windows.  At a given stage choose a
fresh domain shift $s$ and a fresh range shift $t$ so large that a finite
increasing prescription beyond the preceding windows can be extended to an
increasing injection.  Let $\tau_t(n)=t+n$.  Since $d(\tau_t)<\infty$, choose
an increasing injection $\alpha$ and a unit vector $x\in\czz$ such that
\[
  \|J_\alpha\omega(x)-\omega(J_\alpha x)\|_2
   >d(\tau_s)+d(\tau_t)+j
\]
at the $j$th stage.  On $\supp x$ prescribe
\[
 \beta(s+n)=t+\alpha(n).
\]
The prescribed source and target values have the same order.  By taking the
two new windows after the preceding ones and filling the gaps monotonically,
all the finite prescriptions extend to one increasing injection $\beta$.

Exactness makes $\omega$ support-preserving, so the prescription applies both
to $x$ and to $\omega(x)$.  Writing $J_s e_n=e_{s+n}$, on the shifted witness
we therefore have
\begin{align*}
 J_\beta\omega(J_sx)-\omega(J_\beta J_sx)
 &=J_\beta\big(\omega(J_sx)-J_s\omega(x)\big)\\
 &\quad+J_t\big(J_\alpha\omega(x)-\omega(J_\alpha x)\big)\\
 &\quad+\big(J_t\omega(J_\alpha x)-
                  \omega(J_tJ_\alpha x)\big).
\end{align*}
The first and third terms have norms at most $d(\tau_s)$ and $d(\tau_t)$,
while the middle one has the norm of the chosen bad commutator.  Hence the
commutator for the single injection $\beta$ has norm greater than $j$ on a
unit vector for every $j$.  This
contradicts $d(\beta)<\infty$ and proves the lemma.
\end{proof}

\begin{lemma}[Uniform canonical block bound]
\label{lem:uniform-canonical}
There is one constant $K$ for which \eqref{eq:uniform-canonical} holds for
every normalized block basis.
\end{lemma}

\begin{proof}
Let $\kappa(U)$ be the least constant in \eqref{eq:canonical}.  Suppose that
$\sup_U\kappa(U)=\infty$, and let $D$ be given by
Lemma~\ref{lem:uniform-injections}.  Choose normalized block bases $U^j$ and
unit finite-support witnesses $x^j$ for which the left side of
\eqref{eq:canonical} is larger than $j+3D$.

Move the finitely many blocks of $U^j$ used by $x^j$ into a fresh range
window, and move the corresponding coefficient coordinates into a fresh
domain window.  Include the intervening blocks of $U^j$ and fill all other
gaps with normalized singleton blocks.  Doing this successively produces a
single normalized block basis $V=(v_n)$.  If $J_s$ and $J_t$ are the
increasing domain and range embeddings used for one such window, then on the
relevant finite coordinates
\[
 T_VJ_s=J_tT_{U^j}.
\]
Put
\[
 K_U(x)=T_U\omega(x)-\omega(T_Ux)+\sum_nx_n\omega(u_n).
\]
Using the displayed intertwining identity and adding and subtracting
$J_tK_{U^j}(x^j)$ gives
\begin{align*}
 K_V(J_sx^j)-J_tK_{U^j}(x^j)
 &=-T_V\big(J_s\omega(x^j)-\omega(J_sx^j)\big)\\
 &\quad+\big(J_t\omega(T_{U^j}x^j)
                   -\omega(J_tT_{U^j}x^j)\big)\\
 &\quad-\sum_nx_n^j
       \big(J_t\omega(u_n^j)-\omega(J_tu_n^j)\big).
\end{align*}
The first two terms have norm at most $D$.  The vectors in the last sum have
disjoint supports, each has norm at most $D$, and hence that sum also has norm
at most $D\|x^j\|_2=D$.  Thus
\[
 \|K_V(J_sx^j)\|_2>j.
\]
This makes the canonical defect for the single block basis $V$ unbounded,
contradicting Lemma~\ref{lem:canonical}.
\end{proof}

\begin{proof}[Proof of Theorem~\ref{thm:uniform-block}]
The vector estimate is Lemma~\ref{lem:uniform-canonical}.  For the scalar
estimate, by homogeneity it is enough to consider $\sum_np_n=1$.  Put
$x_n=\sqrt{p_n}$ and place $u_n=(\sqrt{q_{n,k}})_k$ in successive disjoint
blocks.  Then $T_Ux=(\sqrt{p_nq_{n,k}})_{n,k}$.  Pair
\eqref{eq:uniform-canonical} with this unit vector.  Disjointness gives
\begin{align*}
 \langle T_U\omega(x),T_Ux\rangle
   &=\mathcal H_\omega(p),\\
 \langle\omega(T_Ux),T_Ux\rangle
   &=\mathcal H_\omega(r),\\
 \left\langle\sum_nx_n\omega(u_n),T_Ux\right\rangle
   &=\sum_np_n\mathcal H_\omega(q^{(n)}).
\end{align*}
Cauchy--Schwarz yields \eqref{eq:entropy-chain}; homogeneity restores the
factor $\sum_np_n$.  In the complex case the final exactification in
Lemma~\ref{lem:full-exact} is a bounded support-preserving perturbation.
Its three terms in \eqref{eq:uniform-canonical} are uniformly bounded (the
last one by disjointness), so the same conclusions hold for the fully exact
representative after enlarging $K$.
\end{proof}

\begin{lemma}[Each multiplier has a bounded commutator]\label{lem:pointwise}
For every $a\in\ell_\infty$, the homogeneous map
\[
 C_a(x)=a\omega(x)-\omega(ax)
\]
is bounded from $(\czz,\|\cdot\|_2)$ to $\ell_2$.
\end{lemma}

\begin{proof}
By homogeneity it suffices to take $\|a\|_\infty\le1$.  Put
$b_n=(1-|a_n|^2)^{1/2}$ and form the normalized block basis
\[
 u_n=a_ne_{2n-1}+b_ne_{2n}.
\]
Let $Jx=\sum_nx_ne_{2n-1}$ and $Kx=\sum_nx_ne_{2n}$, so that
$T_Ux=J(ax)+K(bx)$.  Apply \eqref{eq:canonical} and then project onto the
odd coordinates.  The singleton block bases $(e_{2n-1})$ and $(e_{2n})$
are also covered by \eqref{eq:canonical}; their canonical corrections
vanish because $\omega(e_n)=0$.  Consequently
\[
 J^{-1}P_{\rm odd}\omega(Jz)=\omega(z)+O(\|z\|_2),
 \qquad P_{\rm odd}\omega(Kz)=0,
\]
where the second identity uses \eqref{eq:support}.

Quasilinearity applied to $J(ax)+K(bx)$ now turns the odd-coordinate
projection of \eqref{eq:canonical} into
\begin{equation}\label{eq:multiplier-plus-diagonal}
 \|a\omega(x)-\omega(ax)+d_ax\|_2\le C_a'\|x\|_2,
\end{equation}
where $d_a$ is the diagonal sequence obtained from the odd coordinates of
$\omega(u_n)$.

This diagonal correction is bounded.  Indeed, by quasilinearity,
sign-exactness, and $\omega(e_j)=0$,
\[
 \|\omega(u_n)\|_2
 \le Q(|a_n|+b_n)\le\sqrt2Q.
\]
Thus $\|d_a\|_\infty\le\sqrt2Q$, and
\eqref{eq:multiplier-plus-diagonal} proves that $C_a$ is bounded.
\end{proof}

\begin{lemma}[Uniform multiplier bound]\label{lem:uniform-multiplier}
There is $C<\infty$ such that
\[
 \|C_a(x)\|_2\le C\|a\|_\infty\|x\|_2
 \qquad(a\in\ell_\infty,x\in\czz).
\]
\end{lemma}

\begin{proof}
Set
\[
 q(a)=\sup_{0\ne x\in\czz}\frac{\|C_a(x)\|_2}{\|x\|_2}.
\]
Lemma~\ref{lem:pointwise} says that $q(a)<\infty$ for every
$a\in\ell_\infty$.  Homogeneity and quasilinearity give
\begin{equation}\label{eq:qtriangle}
 q(a+c)\le q(a)+q(c)+Q(\|a\|_\infty+\|c\|_\infty).
\end{equation}

Let $B$ be the closed unit ball of $\ell_\infty$ and
$F_m=\{a\in B:q(a)\le m\}$.  These sets need not be closed, but their
closures have a controlled loss:
\begin{equation}\label{eq:closure-loss}
 \overline{F_m}\subset\{a\in B:q(a)\le m+Q\}.
\end{equation}
To prove this, let $a_j\to a$ in $\ell_\infty$ with $q(a_j)\le m$ and fix
$x\in\czz$.  On each fixed finite-dimensional coordinate subspace,
$\omega$ is bounded: repeated use of quasilinearity and $\omega(e_n)=0$
gives $\|\omega(z)\|_2\le K_{\supp x}\|z\|_2$.  Hence
$\omega((a-a_j)x)\to0$.  Quasilinearity applied to
$ax=a_jx+(a-a_j)x$ then yields
\[
 \|C_a(x)\|_2\le(m+Q)\|x\|_2
\]
after taking the limit.  This proves \eqref{eq:closure-loss}.

Since $B=\bigcup_mF_m\subset\bigcup_m\overline{F_m}$, Baire category and
\eqref{eq:closure-loss} give a norm ball $B(a_0,r)$ contained in $B$ on
which $q$ is bounded by some $M$.  We may choose $a_0$ in the norm interior
of $B$.  If $\|h\|_\infty<r$, then both $a_0$ and $a_0+h$ lie in a slightly
smaller such ball.  Applying \eqref{eq:qtriangle} to
$h=(a_0+h)-a_0$ gives a uniform bound for $q(h)$ on a neighborhood of zero.
Homogeneity now gives $q(a)\le C\|a\|_\infty$ for all $a$.
\end{proof}

\begin{lemma}[Complex exactification]\label{lem:full-exact}
In the complex case, a sign-exact $\ell_\infty$-centralizer $\omega$ is
strongly equivalent to an exact $\ell_\infty$-centralizer
$\widetilde\omega$ satisfying $\widetilde\omega(e_n)=0$.
\end{lemma}

\begin{proof}
Let $K=\TT^{\mathbb N}$.  The centralizer estimate gives, uniformly for
$u\in K$,
\[
 \|\Dop u^{-1}\omega(\Dop u x)-\omega(x)\|_2
 \le C\|x\|_2.
\]
Thus the weak invariant mean
\[
 \widetilde\omega(x)
 =\mathfrak m_{u\in K}\big[\Dop u^{-1}\omega(\Dop u x)\big]
\]
is well-defined in the reflexive space $\ell_2$, differs from $\omega$ by a
bounded homogeneous map, and is exact by invariance of the mean.  It remains
an $\ell_\infty$-centralizer because bounded perturbations preserve the
centralizer estimate.  The equality on the basis is immediate.
\end{proof}

\begin{proof}[Proof of Theorem~\ref{thm:main}]
Normalize the original representative on the basis, apply
Lemmas~\ref{lem:uniform-diagonal} and \ref{lem:exactification}, and retain
self-similarity because equivalent quasilinear maps represent equivalent
exact sequences.  Lemmas~\ref{lem:canonical}--\ref{lem:uniform-multiplier}
then show that the sign-exact representative $\omega$ satisfies the uniform
$\ell_\infty$-centralizer estimate.  Over $\mathbb R$ this is already the
assertion.  Over $\mathbb C$, apply Lemma~\ref{lem:full-exact}.
\end{proof}

\section*{What remains open}

The two theorems remove one layer of generality from the source conjecture:
every self-similar twisted Hilbert extension is represented by an exact
centralizer with a uniform block-natural defect, and its scalar indicator
satisfies the dimension-free ordered chain rule
\eqref{eq:entropy-chain}.  They do \emph{not} prove that this centralizer is
symmetric under coordinate permutations, nor that it has the scalar form
\[
 x_n\varphi\!\left(\log\frac{\|x\|_2}{|x_n|}\right)
\]
treated in \cite{CDR}.  The remaining classification problem can therefore
be stated sharply as follows:

\begin{quote}
Classify the exact $\ell_\infty$-centralizers on $\ell_2$ whose induced
extension is self-similar with respect to every normalized block basis.
Must every nontrivial one generate a space isomorphic to $Z_2$?
\end{quote}

The tempting next step ``self-similar centralizer implies symmetric
centralizer'' is not established here.  Equivalently at the indicator level,
one would like a stable characterization of Shannon entropy from the ordered
refinement rule \eqref{eq:entropy-chain} without assuming permutation
symmetry.  Standard entropy-stability results impose approximate symmetry,
whereas block bases preserve the order of their disjoint supports.  No valid
uniform permutation or ordered-cocycle classification argument was found.

\section*{Verification notes}

The proof was checked against the following failure modes.
\begin{itemize}
\item The first Baire argument does not assume a priori uniform norms for
the block witnesses; the closed sets $E_m$ and their product law supply
that uniformity.
\item The linear corrections may be unbounded on $\ell_2$.  They are only
averaged on $\czz$, where the uniform column bound makes every orbit
bounded; the bounded nonlinear errors are averaged separately.
\item Sign-exactness, rather than an unjustified coordinate formula, supplies
the support property and the blockwise diagonalization in
Lemma~\ref{lem:canonical}.
\item The uniform block constant is not obtained by transporting blocks with
an arbitrary permutation.  Lemmas~\ref{lem:uniform-injections} and
\ref{lem:uniform-canonical} use only increasing domain and range embeddings
and paste finite witnesses into one legitimate block basis.
\item The second Baire argument does not claim that the sublevel sets of
$q$ are closed.  The controlled closure loss \eqref{eq:closure-loss} is
the point that makes the argument valid without continuity of $\omega$ on
the whole unit sphere.
\end{itemize}

No computational verification is used.  Bounded searches on July 21--22, 2026
of the run's 2000--2026 local corpus and web/arXiv results for the exact
conjecture, ``self-similar twisted Hilbert'', ``block centralizer'', and
arXiv:2408.07827 found the source paper and background papers on centralizers
and group actions, but no later resolution or explicit statement of the
reduction or the uniform ordered chain rule above.  Novelty confidence is
therefore moderate, while the
mathematical verdict is \emph{likely valid, expert review requested}.

\begin{thebibliography}{9}

\bibitem{CDR}
W. Corr\^ea, S. Dantas, and D. L. Rodr\'iguez-Vidanes,
\emph{Twisted Hilbert spaces defined by bi-Lipschitz maps},
arXiv:2408.07827 (2024).

\bibitem{CG}
F. Cabello S\'anchez and R. Garc\'ia,
\emph{The twisted Hilbert space ideals},
Integral Equations Operator Theory 94 (2022), Paper No.~21;
arXiv:2112.03852.

\bibitem{CF}
J. M. F. Castillo and V. Ferenczi,
\emph{Group actions on twisted sums of Banach spaces},
Bull. Malays. Math. Sci. Soc. 46 (2023), Paper No.~135;
arXiv:2003.09767.

\end{thebibliography}

\end{document}
